\lyseis{Β1.6}{Είδη γωνιών και κάθετες}
\sol{1} $37^\circ$ οξεία, $90^\circ$ ορθή, $123^\circ$ αμβλεία, $180^\circ$ ευθεία, $250^\circ$ μη κυρτή, $360^\circ$ πλήρης, $0^\circ$ μηδενική, $89^\circ$ οξεία.
\sol{2} Ελεύθερη κατασκευή. Οι διχοτόμοι σχηματίζουν το μισό της γωνίας.
\sol{3} $2{:}00$: $60^\circ$ οξεία. $3{:}00$: $90^\circ$ ορθή. $4{:}00$: $120^\circ$ αμβλεία. $6{:}00$: $180^\circ$ ευθεία. $7{:}00$: $150^\circ$ αμβλεία (μικρότερη γωνία).
\sol{4} Οι άλλες τρεις είναι ορθές.
\sol{5} Ελεύθερη κατασκευή. Από κάθε σημείο υπάρχει μία μόνο κάθετη.
\sol{6} Μία κάθετη και άπειρες πλάγιες.
\sol{7} α) Σωστό. β) Σωστό. γ) Λάθος: είναι όλες ορθές. δ) Λάθος: η μηδενική δεν είναι οξεία. ε) Λάθος: μη κυρτή.
\sol{8} Διπλώνουμε το φύλλο και ταιριάζουμε μια γωνία του φύλλου με τη γωνία του φράχτη. Αν εφαρμόζει, είναι ορθή.
\sol{9} Ορθή γωνία ανάμεσα στις πλευρές $3$ και $4$.
\sol{10} Ελεύθερη απάντηση. Π.χ. ορθές: γωνίες του πίνακα, της πόρτας, του θρανίου.
\sol{11} Η $110^\circ$ είναι αμβλεία. Μια ορθή γωνία είναι $90^\circ$.
\sol{12} Προς Ανατολή: $90^\circ$. Προς Νότο: $180^\circ$.
\sol{13} Η ορθή και η διπλανή της αθροίζουν $180^\circ$, άρα η διπλανή είναι $90^\circ$. Η απέναντι είναι ίση με την αρχική και η τέταρτη ίση με τη διπλανή, άρα $90^\circ$.
\sol{14} Όχι: δύο ορθές θα είχαν άθροισμα $180^\circ$ και δεν θα έμενε μέτρο για την τρίτη γωνία.
\sol{15} Ο γνώμονας έχει μία ορθή γωνία. Για να τον ελέγξεις, φέρε κάθετη και κατόπιν ανάποδα· αν ταιριάζουν, είναι σωστός.
\sol{16} Η $180^\circ-\omega$ είναι αμβλεία (μεταξύ $90^\circ$ και $180^\circ$).
\sol{17} Ελεύθερη απάντηση.
\sol{18} Ελεύθερη απάντηση.
